% ECONOMICS_PROBLEM: {"id": "D83-218", "title": "Approximating receiver-optimal cheap-talk equilibria", "jel_code": "D83", "source_id": "OP-0178"}
% !TeX program = xelatex
\documentclass[11pt,a4paper]{article}
\usepackage[margin=23mm]{geometry}
\usepackage{fontspec}
\setmainfont{Latin Modern Roman}
\usepackage{amsmath,amssymb,mathtools}
\usepackage{xcolor,enumitem,booktabs,longtable,array,xurl,hyperref}
\definecolor{muted}{HTML}{53636C}
\hypersetup{colorlinks=true,urlcolor=blue,linkcolor=blue}
\setlength{\parindent}{0pt}
\setlength{\parskip}{5pt}
\newcommand{\R}{\mathbb R}
\newcommand{\E}{\mathbb E}
\newcommand{\N}{\mathbb N}
\newcommand{\ind}{\mathbf 1}
\newcommand{\Var}{\operatorname{Var}}
\newcommand{\Cov}{\operatorname{Cov}}
\newcommand{\supp}{\operatorname{supp}}
\DeclareMathOperator*{\argmin}{arg\,min}
\DeclareMathOperator*{\argmax}{arg\,max}
\newcommand{\entrymeta}[3]{{\sffamily\small\color{muted}\textbf{\texttt{#1}}\quad|\quad #2\par #3\par}}
\newcommand{\Problemref}[1]{\texttt{#1}}
\newcommand{\JELref}[2]{\texttt{#2} (JEL \texttt{#1})}
\begin{document}
\section*{Approximating receiver-optimal cheap-talk equilibria}
\textbf{Problem ID:} \texttt{D83-218}\quad \textbf{JEL:} \texttt{D83}\par
% BEGIN ECONOMICS STATEMENT
\label{problem:OP-0178}
% GAME_THEORY_PROBLEM: {"local_id": "GT-048", "original_number": 48, "kind": "R", "entry_kind": "statement_only", "published": false, "review_required": true, "category": "game_theory"}
\label{game:GT-048}
{\sffamily\small\color{muted}
\textbf{GT-048} | Strategic communication and complexity | \textbf{R}: normalized approximation formulation of a source agenda.
\par}
\subsection*{Definitions and mathematical statement}
An input consists of finite $\Omega,A$, a rational prior $\mu$ with full support, and rational payoffs $u_S,u_R:\Omega\times A\to[0,1]$. Fix a message set $M$ of size $|\Omega|+1$. A strategy pair is $\pi(m\mid\omega)$ and $r(a\mid m)$. Put $q_m=\sum_\omega\mu(\omega)\pi(m\mid\omega)$. A cheap-talk equilibrium requires, for all $\omega,m,m'$,
\[
 \pi(m\mid\omega)>0\ \Longrightarrow\
 \sum_a r(a\mid m)u_S(\omega,a)\geq\sum_a r(a\mid m')u_S(\omega,a),
\]
and, whenever $q_m>0$ and $r(a\mid m)>0$,
\[
 \sum_\omega\mu(\omega)\pi(m\mid\omega)u_R(\omega,a)
 \geq\sum_\omega\mu(\omega)\pi(m\mid\omega)u_R(\omega,b)
 \quad\forall b\in A.
\]
For unused messages require receiver optimality under some belief in $\Delta(\Omega)$. Let $U_R(\pi,r)=\sum_{\omega,m,a}\mu(\omega)\pi(m\mid\omega)r(a\mid m)u_R(\omega,a)$ and $U_R^*=\sup_{\rm equilibria}U_R$. Characterize constants $c\in(0,1]$ for which a polynomial-time algorithm returns an exact equilibrium with $U_R\geq cU_R^*$, or gives an additive guarantee $U_R\geq U_R^*-\varepsilon$ for each fixed $\varepsilon>0$. Exact algebraic probabilities must be output in an explicit algebraic-number representation.
\subsection*{Short English statement}
Determine how well receiver welfare can be approximated while respecting both parties' cheap-talk incentives.
\subsection*{Variants and scope boundaries}
Approximate equilibrium and exact equilibrium are distinct outputs. Multiplicative approximation requires a fixed nonnegative utility normalization. Sender commitment would change cheap talk into persuasion.
\subsection*{Source relationship and status}
The source proves receiver threshold hardness and discusses the absence of positive results for receiver utility. The explicit $c$-approximation family here is editorial; it is not attributed as a verbatim conjecture. Proposition 3.10 justifies the bounded-message optimum.
\subsection*{References}
\begingroup\footnotesize\raggedright
\textbf{[S1]} Yakov Babichenko, Inbal Talgam-Cohen, Haifeng Xu, and Konstantin Zabarnyi (2024 version; first posted 2023). \emph{Algorithmic Cheap Talk}. arXiv:2311.09011v2. \textit{Verified locator:} Section 2; Section 3.2, Propositions 3.10--3.11; Section 5, ``Maximization objective.'' \url{https://arxiv.org/html/2311.09011v2}
\par\endgroup

\subsection*{Common conventions: Source mathematical conventions}

\textbf{Finite games.} For a finite set $A$, $\Delta(A)$ is its probability
simplex. Players choose independent mixed strategies in Nash equilibrium;
correlation is allowed only when explicitly part of the solution concept.
In a game with payoff functions $u_i$ and strategy profile $x$, an additive
$\varepsilon$ Nash equilibrium satisfies
\[
 u_i(x)\geq u_i(x_i',x_{-i})-\varepsilon
 \quad\text{for every player $i$ and every admissible deviation $x_i'$}.
\]
For numerical approximation, payoffs are normalized as locally specified.
An $\varepsilon$ payoff residual is distinct from distance at most
$\varepsilon$ to the set of exact equilibria. Point convergence, convergence
of empirical play, and convergence in distance to an equilibrium set are
also distinct.

\textbf{Computational representations.} Unless stated otherwise, a complexity
question takes rational data with binary encoding; polynomial time is measured
in total bit length. A fixed-accuracy polynomial algorithm is not necessarily
polynomial in $1/\varepsilon$, and neither is necessarily polynomial in
$\log(1/\varepsilon)$. Succinct games and value, demand, or sampling oracles
must declare their interfaces. Exact real solutions require an explicit
representation; an unspecified list of arbitrary real numbers is not a
finite computational output. A claimed algorithmic barrier is meaningful
only for its specified model and reductions.

\textbf{Dynamic games.} Histories, information, observation, and allowable
randomization are part of the model. A stationary strategy depends only on
the current state; a positional strategy in a deterministic graph game
chooses one action at each controlled vertex. Nash equilibrium at an initial
state and equilibrium after every history are different requirements.
An expectation of a pathwise $\liminf$ payoff, a $\liminf$ of expected averages,
a limiting expected average, and a uniform finite-horizon equilibrium are
different criteria. None is silently substituted for another.

\textbf{Allocation and mechanisms.} A complete allocation assigns every item
exactly once unless otherwise stated. Zero-valued goods, negative values,
capacity constraints, feasibility, lotteries, and copies of goods affect
fairness definitions and are specified locally. Dominant-strategy incentive
compatibility, Bayesian incentive compatibility, universal truthfulness,
and truthfulness in expectation impose different inequalities. Whenever
an objective ratio has zero denominator, its convention is stated or those
instances are excluded.

\textbf{Infinite-dimensional models.} Expectations, integrals, stochastic
equations, and optimization objectives require their local regularity,
integrability, filtrations, and admissible domains. A backward--forward
mean-field system is a boundary-value problem; it is not an initial-value
semiflow without an additional construction. Long-time, large-population,
and vanishing-noise limits require an order of limits and a convergence
topology. If a minimum need not be attained, the objective is its infimum
or supremum and the approximation requirement is stated separately.
% END ECONOMICS STATEMENT
\end{document}
